\DocumentMetadata{
  pdfversion=2.0,
  pdfstandard=ua-2,
  lang=en-US,
  tagging=on,
  tagging-setup={math/setup=mathml-SE,tabsorder=structure}
}
\documentclass[11pt,letterpaper]{article}
\usepackage[margin=0.68in,includefoot]{geometry}
\usepackage{newcomputermodern}
\usepackage{amsmath,amscd,mathtools}
\usepackage{tikz,adjustbox}
\providecommand{\square}{\mdlgwhtsquare}
\usepackage[hidelinks]{hyperref}
\hypersetup{
  pdftitle={Analysis PhD Exam, May 2010},
  pdfauthor={University of Florida Department of Mathematics},
  pdfsubject={Graduate mathematics examination},
  pdfdisplaydoctitle=true,
  bookmarksopen=true
}
\setcounter{secnumdepth}{0}
\setlength{\parindent}{0pt}
\setlength{\parskip}{0.28em}
\setlength{\emergencystretch}{3em}
\setlength{\leftmargini}{2.15em}
\setlength{\leftmarginii}{2.65em}
\setlength{\leftmarginiii}{3em}
\pagestyle{empty}
\raggedbottom
\allowdisplaybreaks
\NewTaggingSocketPlug{block/list/label}{examproblem}{%
  \tagstructbegin{tag=Lbl}\tagstructend%
  \tagstructbegin{tag=\LBody}%
  \tagstructbegin{tag=\ProblemHeadingTag,title-o={\ProblemHeadingText}}%
  \tagmcbegin{tag=\ProblemHeadingTag,actualtext={\ProblemHeadingText}}%
  #2%
  \tagmcend%
  \tagstructend%
}
\newenvironment{examproblems}[2]{%
  \def\ProblemHeadingTag{#2}%
  \AssignTaggingSocketPlug{block/list/label}{examproblem}%
  \begin{enumerate}%
  \setcounter{enumi}{#1}%
  \renewcommand{\labelenumi}{\textbf{\theenumi.}}%
  \setlength{\topsep}{0.35em}%
  \setlength{\partopsep}{0pt}%
  \setlength{\itemsep}{0.48em}%
  \setlength{\parsep}{0.18em}%
}{\end{enumerate}}
\newenvironment{examsubparts}{%
  \AssignTaggingSocketPlug{block/list/label}{default}%
  \begin{enumerate}%
  \setlength{\topsep}{0.18em}%
  \setlength{\partopsep}{0pt}%
  \setlength{\itemsep}{0.16em}%
  \setlength{\parsep}{0.1em}%
}{\end{enumerate}}
\newenvironment{examinstructions}{%
  \par\begingroup\itshape
}{%
  \par\endgroup\vspace{0.18em}
}
\makeatletter
\renewcommand\subsection{\@startsection{subsection}{2}{0pt}%
  {-1.25ex plus -.25ex minus -.1ex}%
  {0.55ex plus .1ex}%
  {\normalfont\large\bfseries}}
\renewcommand{\@maketitle}{%
  \newpage\null\vskip -1em%
  {\footnotesize\bfseries
    \href{https://www.ufl.edu/}{University of Florida}, \href{https://math.ufl.edu/}{Department of Mathematics}\par}%
  \vskip -0.5em%
  \tagstructbegin{tag=H1,title={Analysis PhD Exam, May 2010}}%
  \tagpdfparaOff%
  \tagmcbegin{tag=H1}%
    {\LARGE\bfseries\href{https://gma.math.ufl.edu/exams/analysis-phd/}{Analysis PhD Exam}, May 2010\par}%
  \tagmcend%
  \tagpdfparaOn%
  \tagstructend%
  \vskip 0.25em%
  \tagmcbegin{artifact}%
    \hrule height 1pt%
  \tagmcend%
  \vskip 1em%
}
\makeatother
\title{Analysis PhD Exam, May 2010}
\author{}
\date{}
\begin{document}
\maketitle
\thispagestyle{empty}
\begin{examinstructions}
Write solutions in a neat and logical fashion, giving complete reasons for all steps. Do any seven problems.
\end{examinstructions}
\begin{examproblems}{0}{H2}
\def\ProblemHeadingText{Problem 1.}
\item
Consider Lebesgue measure on the real line~\(\mathbb{R}\). Give an example for each of the following, if possible. If not possible, give a brief explanation. A sequence \(f_n\) in \(L^1(\mathbb{R})\) converging to an~\(f\) in \(L^1(\mathbb{R})\)... Which of these answers change if we replace~\(\mathbb{R}\) by the interval \([0,1]\)?
\begin{examsubparts}
\item[\textnormal{a)}]
...in the \(L^1\) norm but not in measure,
\item[\textnormal{b)}]
...in measure but not in the \(L^1\) norm,
\item[\textnormal{c)}]
...in the \(L^1\) norm but not a.e.,
\item[\textnormal{d)}]
...a.e. but not in measure.
\end{examsubparts}
\def\ProblemHeadingText{Problem 2.}
\item
Define the Hardy–Littlewood maximal function. State and prove the Hardy–Littlewood maximal theorem. (Begin by proving an appropriate covering lemma.)
\def\ProblemHeadingText{Problem 3.}
\item
This problem has two parts.
\begin{examsubparts}
\item[\textnormal{a)}]
State Egorov’s theorem.
\item[\textnormal{b)}]
Suppose \((X,\mathcal{M},\mu)\) is a measure space with \(\mu(X)<\infty\) and \(f:X\to\mathbb{C}\) is measurable. Let \((f_n)\) be a sequence of integrable functions such that \(f_n\to f\) a.e. Suppose further that the sequence \(f_n\) is uniformly integrable, that is, for every \(\epsilon>0\) there exists a \(\delta>0\) such that if~\(E\) is measurable and \(\mu(E)<\delta\), then
\[
\int_E |f_n|\,d\mu<\epsilon.
\]
 Prove that~\(f\) is integrable and
\[
\lim_{n\to\infty}\int_X |f_n-f|\,d\mu=0.
\]
\end{examsubparts}
\def\ProblemHeadingText{Problem 4.}
\item
Let \(\mathcal{X},\mathcal{Y}\) be Banach spaces. Suppose that \((T_n)\) is a sequence of bounded linear operators from~\(\mathcal{X}\) to~\(\mathcal{Y}\) such that \(\lim_{n\to\infty}T_nx\) exists for every \(x\in\mathcal{X}\). Prove that
\[
Tx:=\lim_{n\to\infty}T_nx
\]
 defines a bounded linear operator from~\(\mathcal{X}\) to~\(\mathcal{Y}\).
\def\ProblemHeadingText{Problem 5.}
\item
This problem has two parts.
\begin{examsubparts}
\item[\textnormal{a)}]
State the Lebesgue–Radon–Nikodym theorem.
\item[\textnormal{b)}]
State and prove a version of the chain rule for Radon–Nikodym derivatives.
\end{examsubparts}
\def\ProblemHeadingText{Problem 6.}
\item
Let~\(X\) be a normed vector space. Given \(x\in X\), define a linear functional \(\hat{x}:X^*\to\mathbb{C}\) by \(\hat{x}(f)=f(x)\). Prove that the map \(x\to\hat{x}\) is an isometry from~\(X\) into \(X^{**}\).
\def\ProblemHeadingText{Problem 7.}
\item
Let~\(\mathcal{H}\) be a Hilbert space and suppose \(\mathcal{M},\mathcal{N}\) are closed subspaces with \(\mathcal{M}\perp\mathcal{N}\). Prove that \(\mathcal{M}+\mathcal{N}\) is closed.
\def\ProblemHeadingText{Problem 8.}
\item
Suppose \((k_n)_{n=1}^{\infty}\) is a sequence of functions in \(L^1(\mathbb{R})\) satisfying the following. Prove that
\[
\lim_{n\to\infty}\|f-f*k_n\|_1=0
\]
 for all \(f\in L^1(\mathbb{R})\). (Here \(*\) denotes convolution: \((f*g)(x):=\int_{-\infty}^{\infty}f(x-t)g(t)\,dt\).)
\begin{examsubparts}
\item[\textnormal{a)}]
\(\displaystyle \int_{-\infty}^{\infty}k_n(t)\,dt=1\) for all~\(n\),
\item[\textnormal{b)}]
\(\displaystyle \sup_n\int_{-\infty}^{\infty}|k_n(t)|\,dt<\infty\), and
\item[\textnormal{c)}]
for each \(\delta>0\), \(\displaystyle \sup_{|t|\geq\delta}\{|k_n(t)|\}\to0\) as \(n\to\infty\).
\end{examsubparts}
\end{examproblems}
\end{document}
